\chapter{Fasting for a Blood Test}
\section*{What Is This Test?} Fasting is a preparation requirement before selected blood tests, usually meaning no food or caloric drinks for a specified period.
\section*{Alternative Names} Fasting blood work; fasting laboratory test preparation.
\section*{Principle} Avoiding recent food intake reduces short-term effects on glucose, triglycerides, and some chemistry results.
\section*{Why This Test Is Ordered} Fasting may be requested for lipid testing, glucose testing, or specific metabolic measurements, although many tests do not require it.
\section*{Primary vs Secondary Test} Fasting is preparation, not a diagnostic test; the ordered blood test supplies the clinical result.
\section*{How This Test Is Performed} The patient follows the instructed fast, provides a blood sample, and resumes normal intake afterward unless told otherwise.
\section*{What Preparation Is Needed} Ask how long to fast and whether water, medicines, smoking, or exercise are allowed. Do not fast or alter diabetes medicines without instructions.
\section*{Understanding Results} Nonfasting collection can affect tests differently; record fasting duration so the laboratory and clinician can interpret results.
\section*{Follow-up Tests} Repeat fasting collection or a nonfasting confirmatory test may be recommended if preparation was inadequate.
\section*{References} \begin{enumerate}\item MedlinePlus. Fasting for a Blood Test.\item Clinical and Laboratory Standards Institute. Patient preparation guidance.\end{enumerate}
\clearpage\section*{Understanding Results and Follow-up}
The laboratory report should identify the specimen, method, units, reference interval or decision limit, and whether the result is positive, negative, abnormal, or inconclusive. Reference intervals vary with age, sex, pregnancy, specimen type, assay, and laboratory; a result outside a range is not automatically a diagnosis, and a result within a range does not always exclude disease. Discuss unexpected results with the ordering clinician, who can decide whether repeat or confirmatory testing, treatment, or referral is appropriate. Do not change medicines or delay urgent care based on an isolated result.
\section*{Regulatory Status and Authorization Date}This chapter describes a test category, not a specific commercial product. FDA, CDSCO, EU IVDR, and MHRA approval or registration dates therefore cannot be assigned without the assay manufacturer, product name, instrument, and intended use. For laboratory-developed tests, an individual agency approval date may not exist. See the Preface for the interpretation of regulatory status.
\clearpage
